% Options for packages loaded elsewhere
\PassOptionsToPackage{unicode}{hyperref}
\PassOptionsToPackage{hyphens}{url}
\PassOptionsToPackage{dvipsnames,svgnames,x11names}{xcolor}
%
\documentclass[
  11pt,
]{article}
\usepackage{amsmath,amssymb}
\usepackage{iftex}
\ifPDFTeX
  \usepackage[T1]{fontenc}
  \usepackage[utf8]{inputenc}
  \usepackage{textcomp} % provide euro and other symbols
\else % if luatex or xetex
  \usepackage{unicode-math} % this also loads fontspec
  \defaultfontfeatures{Scale=MatchLowercase}
  \defaultfontfeatures[\rmfamily]{Ligatures=TeX,Scale=1}
\fi
\usepackage{lmodern}
\ifPDFTeX\else
  % xetex/luatex font selection
\fi
% Use upquote if available, for straight quotes in verbatim environments
\IfFileExists{upquote.sty}{\usepackage{upquote}}{}
\IfFileExists{microtype.sty}{% use microtype if available
  \usepackage[]{microtype}
  \UseMicrotypeSet[protrusion]{basicmath} % disable protrusion for tt fonts
}{}
\makeatletter
\@ifundefined{KOMAClassName}{% if non-KOMA class
  \IfFileExists{parskip.sty}{%
    \usepackage{parskip}
  }{% else
    \setlength{\parindent}{0pt}
    \setlength{\parskip}{6pt plus 2pt minus 1pt}}
}{% if KOMA class
  \KOMAoptions{parskip=half}}
\makeatother
\usepackage{xcolor}
\usepackage[margin=1in]{geometry}
\usepackage{longtable,booktabs,array}
\usepackage{calc} % for calculating minipage widths
% Correct order of tables after \paragraph or \subparagraph
\usepackage{etoolbox}
\makeatletter
\patchcmd\longtable{\par}{\if@noskipsec\mbox{}\fi\par}{}{}
\makeatother
% Allow footnotes in longtable head/foot
\IfFileExists{footnotehyper.sty}{\usepackage{footnotehyper}}{\usepackage{footnote}}
\makesavenoteenv{longtable}
\setlength{\emergencystretch}{3em} % prevent overfull lines
\providecommand{\tightlist}{%
  \setlength{\itemsep}{0pt}\setlength{\parskip}{0pt}}
\setcounter{secnumdepth}{-\maxdimen} % remove section numbering
\ifLuaTeX
  \usepackage{selnolig}  % disable illegal ligatures
\fi
\IfFileExists{bookmark.sty}{\usepackage{bookmark}}{\usepackage{hyperref}}
\IfFileExists{xurl.sty}{\usepackage{xurl}}{} % add URL line breaks if available
\urlstyle{same}
\hypersetup{
  pdfauthor={Stefano Alimonti --- stefalimonti@gmail.com},
  colorlinks=true,
  linkcolor={blue},
  filecolor={Maroon},
  citecolor={Blue},
  urlcolor={blue},
  pdfcreator={LaTeX via pandoc}}

\author{Stefano Alimonti --- stefalimonti@gmail.com}
\date{March 2026}

\begin{document}

{
\hypersetup{linkcolor=black}
\setcounter{tocdepth}{2}
\tableofcontents
}
\section{The Carry-Zero Entropy Bound: Structural Limits of Bitwise
Factorization}\label{the-carry-zero-entropy-bound-structural-limits-of-bitwise-factorization}

\textbf{Author:} Stefano Alimonti \textbf{Affiliation:} Independent Researcher
\textbf{Date:} March 2026

\begin{center}\rule{0.5\linewidth}{0.5pt}\end{center}

\newpage

\subsection{Abstract}\label{abstract}

The carry polynomial framework represents positional multiplication
\(N = p \cdot q\) as the polynomial relation \(P(x)Q(x) = F_N(x) + (x-2)C(x)\).
We analyze the 2-adic tree of partial factorizations generated by bit-by-bit BFS
in base 2. We define the \emph{carry-zero envelope} \(S_0(k)\) as the number of
partial-factor pairs \((p_k, q_k)\) at bit position \(k\) with zero carry-out,
and prove a structural decomposition:

\[S_0(k) = \tau_{\text{odd}}(M_k) + \Phi(k)\]

where \(\tau_{\text{odd}}(M_k)\) counts exact odd-divisor pairs of
\(M_k = N \bmod 2^{k+1}\) and \(\Phi(k)\) counts \emph{phantom} states --- pairs
satisfying the local carry constraint whose integer product exceeds \(M_k\). We
prove that \(\tau_{\text{odd}}(M_k)\) is sub-polynomially bounded via
Wigert--Ramanujan, and present empirical evidence that \(\Phi(k)\) grows
exponentially as \(\sim 2^{\alpha D}\) with \(\alpha \in [0.38, 0.50]\),
dominating \(S_0(k)\) for large \(D\). This reveals a clean structural
separation: a vanishing \emph{core} of real factorizations embedded in an
exponential \emph{halo} of \(p\)-adic modular aliases.

\textbf{Keywords:} integer factorization, carry propagation, 2-adic numbers,
Hensel lifting, divisor function, Wigert-Ramanujan

\textbf{MSC 2020:} 11A63, 11Y05, 11S85, 68Q25

\begin{center}\rule{0.5\linewidth}{0.5pt}\end{center}

\newpage

\subsection{1. Introduction}\label{introduction}

The factorization of a semiprime \(N = p \cdot q\) of bit-length \(D\) is
computationally hard. The best general-purpose algorithms (GNFS) achieve
sub-exponential but super-polynomial complexity. Simpler approaches operate in
\(O(\sqrt{N}) = O(2^{D/2})\).

The Carry Representation Theorem {[}B{]} establishes that binary multiplication
satisfies \(P(x)Q(x) = F_N(x) + (x-2)C(x)\), where \(P\), \(Q\), \(F_N\) are
digit polynomials and \(C\) is the carry polynomial. A natural question: can the
carry structure be exploited to factor \(N\) more efficiently?

We answer negatively by analyzing the BFS tree of partial factorizations
constrained by \(N \bmod 2^k\). The carry-zero envelope --- states where no
information overflows to the next bit --- decomposes into a sub-polynomial core
(real divisors) and an exponential halo (modular phantoms). The halo arises from
the degeneracy of 2-adic Hensel lifting and provides exponential camouflage for
the true factors.

\emph{Scope remark.} This paper is a focused companion to {[}B{]}, which
establishes the carry polynomial framework. We prove the structural
decomposition (Theorem 1) and the sub-polynomial core bound (Theorem 2), and
present empirical evidence for the phantom growth rate. The paper is
deliberately concise; readers seeking the broader algebraic context (CRT,
companion matrices, spectral determinants) are directed to {[}B{]}.

\begin{center}\rule{0.5\linewidth}{0.5pt}\end{center}

\newpage

\subsection{2. The Carry-Zero Envelope}\label{the-carry-zero-envelope}

\textbf{Definition 1} (BFS state). At step \(k\), a state is a triple
\((p_k, q_k, c_k)\) where \(p_k, q_k\) are \((k+1)\)-bit odd numbers and
\(c_k \in \mathbb{Z}_{\geq 0}\) is the carry into column \(k+1\):

\[c_{k+1} = \left\lfloor \frac{conv_k + c_k}{2} \right\rfloor, \qquad b_k(N) = (conv_k + c_k) \bmod 2\]

where \(conv_k = \sum_{i+j=k} b_i(p_k) \cdot b_j(q_k)\).

\textbf{Definition 2} (Carry-zero envelope).
\(S_0(k) = |\{(p_k, q_k) : c_{k+1} = 0\}|\).

\subsubsection{2.1 Local vs Global Carry}\label{local-vs-global-carry}

A carry of zero at column \(k\) means \(conv_k + c_k < 2\): no overflow into
column \(k+1\). This is a \emph{local} constraint on a single column. It does
\textbf{not} imply \(p_k \cdot q_k = M_k\) as integers, where
\(M_k = N \bmod 2^{k+1}\).

\textbf{Example.} For \(N = 15\), at \(k = 2\): \(M_2 = 7\). The pair \((5, 3)\)
has carry \(= 0\) at columns 0--2, yet \(5 \times 3 = 15 \neq 7\). The lower 3
bits match (\(15 \bmod 8 = 7\)), but the product overflows via higher-order
convolution terms.

\textbf{Implementation note.} In the BFS enumerator (D03), the carry-zero
condition \(c_{k+1} = 0\) is tested via the equivalent criterion \(P_k = M_k\),
where

\[P_k = \sum_{j=0}^{k} \mathrm{conv}_j \cdot 2^j\]

is the carry-free convolution sum. Since BFS enforces
\((p_k \cdot q_k) \bmod 2^{k+1} = M_k\), the identity \(P_k = M_k\) holds if and
only if no carry propagation occurs at columns \(0, \ldots, k\).

\begin{center}\rule{0.5\linewidth}{0.5pt}\end{center}

\newpage

\subsection{3. Decomposition Theorem}\label{decomposition-theorem}

\textbf{Theorem 1} (Decomposition).
\(S_0(k) = \tau_{\text{odd}}(M_k) + \Phi(k)\), where:

\begin{itemize}
\tightlist
\item
  \(\tau_{\text{odd}}(M_k) = |\{(d,\, M_k/d) : d \mid M_k,\; d \text{ and } M_k/d \text{ both odd}\}|\)
\item
  \(\Phi(k) = |\{(p_k, q_k) : c_{k+1} = 0,\; p_k \cdot q_k \neq M_k\}|\)
  (phantoms)
\end{itemize}

\emph{Proof.} If \(c_{k+1} = 0\) at all columns \(0, \ldots, k\), then
\(p_k \cdot q_k \equiv M_k \pmod{2^{k+1}}\). Among these, exactly those with
\(p_k \cdot q_k = M_k\) (as integers) correspond to odd-divisor pairs. The
remainder have \(p_k \cdot q_k = M_k + j \cdot 2^{k+1}\) for some \(j \geq 1\).
\(\square\)

\begin{center}\rule{0.5\linewidth}{0.5pt}\end{center}

\newpage

\subsection{4. The Divisor Core: Sub-Polynomial
Bound}\label{the-divisor-core-sub-polynomial-bound}

\textbf{Theorem 2} (Wigert--Ramanujan bound).
\(\max_k \tau_{\text{odd}}(M_k) \leq \max_{X < 2^D} \tau(X) = 2^{O(D / \log_2 D)}\).

In particular,
\(\alpha_\tau := \log_2(\max_k \tau_{\text{odd}}(M_k)) / D \to 0\) as
\(D \to \infty\).

\emph{Proof.}
\(\tau_{\text{odd}}(M_k) \leq \tau(M_k) \leq \max_{X < 2^D} \tau(X)\). The
Wigert--Ramanujan theorem gives
\(\max_{X < Y} \tau(X) = 2^{(1+o(1)) \ln Y / \ln \ln Y}\). Substituting
\(Y = 2^D\): \(\log_2(\max \tau) = O(D / \log_2 D)\). \(\square\)

\begin{center}\rule{0.5\linewidth}{0.5pt}\end{center}

\newpage

\subsection{5. Analytical Bounds on the Phantom
Exponent}\label{analytical-bounds-on-the-phantom-exponent}

\subsubsection{5.0 The 2-adic Solution Count}\label{the-2-adic-solution-count}

\textbf{Proposition 3} (Upper bound). For \(M_k\) odd, define

\[\mathcal{M}(k) := \lvert \{(p,q) : p,q \text{ odd},\; 1 \leq p,q \leq 2^{k+1}-1,\; p \cdot q \equiv M_k \pmod{2^{k+1}} \} \rvert.\]

Then \(\mathcal{M}(k) = 2^k\) exactly. Since \(S_0(k) \leq \mathcal{M}(k)\):

\[\alpha_{S_0} := \frac{\log_2 S_0(k)}{k} \leq 1\]

\emph{Proof.} Since \(M_k\) is odd (as \(N\) is a product of odd primes),
\(M_k \in (\mathbb{Z}/2^{k+1}\mathbb{Z})^{\times}\). For each odd
\(p \in \{1, 3, 5, \ldots, 2^{k+1}-1\}\) (there are \(2^k\) such values), the
equation \(p \cdot q \equiv M_k \pmod{2^{k+1}}\) has exactly one solution
\(q \equiv p^{-1} M_k \pmod{2^{k+1}}\), and \(q\) is odd because \(p^{-1} M_k\)
is odd (product of odd units). Every carry-zero pair \((p_k, q_k)\) satisfies
\(p_k \cdot q_k \equiv M_k \pmod{2^{k+1}}\) (Theorem 1 proof), so
\(S_0(k) \leq \mathcal{M}(k) = 2^k\). \(\square\)

\textbf{Corollary 3.1.} \(\Phi(k) \leq 2^k - \tau_{\text{odd}}(M_k)\). Combined
with Theorem 2: \(\Phi(k) \leq 2^k - 1\), so \(\alpha_\Phi \leq 1\).

The empirical finding \(\alpha \in [0.38, 0.50]\) (below) shows that the
carry-zero constraint at \emph{every} intermediate column eliminates a constant
fraction of the exponent: roughly half of all 2-adic solutions fail the
column-by-column carry test. The gap \(1 - \alpha\) quantifies the cumulative
filtering power of carry propagation constraints.

\subsubsection{5.1 Experimental Results}\label{experimental-results}

Phantom states arise from modular aliasing: pairs whose product matches \(M_k\)
modulo \(2^{k+1}\) but overflows. Unlike \(\tau_{\text{odd}}\), phantoms grow
exponentially.

Measured via exact BFS (Python for \(D \leq 20\); C hash-table BFS for
\(D \leq 28\)) on random balanced semiprimes:

\begin{longtable}[]{@{}
  >{\raggedleft\arraybackslash}p{(\columnwidth - 12\tabcolsep) * \real{0.1429}}
  >{\raggedleft\arraybackslash}p{(\columnwidth - 12\tabcolsep) * \real{0.1429}}
  >{\raggedleft\arraybackslash}p{(\columnwidth - 12\tabcolsep) * \real{0.1429}}
  >{\raggedleft\arraybackslash}p{(\columnwidth - 12\tabcolsep) * \real{0.1429}}
  >{\raggedleft\arraybackslash}p{(\columnwidth - 12\tabcolsep) * \real{0.1429}}
  >{\raggedleft\arraybackslash}p{(\columnwidth - 12\tabcolsep) * \real{0.1429}}
  >{\raggedleft\arraybackslash}p{(\columnwidth - 12\tabcolsep) * \real{0.1429}}@{}}
\toprule\noalign{}
\begin{minipage}[b]{\linewidth}\raggedleft
\(D\)
\end{minipage} & \begin{minipage}[b]{\linewidth}\raggedleft
\(\langle S_0 \rangle\)
\end{minipage} & \begin{minipage}[b]{\linewidth}\raggedleft
\(\langle \tau \rangle\)
\end{minipage} & \begin{minipage}[b]{\linewidth}\raggedleft
\(\langle \Phi \rangle\)
\end{minipage} & \begin{minipage}[b]{\linewidth}\raggedleft
\% Phantom
\end{minipage} & \begin{minipage}[b]{\linewidth}\raggedleft
\(\alpha_{S_0}\)
\end{minipage} & \begin{minipage}[b]{\linewidth}\raggedleft
\(\alpha_{\tau}\)
\end{minipage} \\
\midrule\noalign{}
\endhead
\bottomrule\noalign{}
\endlastfoot
5 & 8 & 4.0 & 4.0 & 50\% & 0.60 & 0.40 \\
8 & 12 & 4.0 & 8.0 & 67\% & 0.45 & 0.25 \\
10 & 34 & 4.0 & 30 & 88\% & 0.51 & 0.20 \\
12 & 56 & 3.7 & 52 & 93\% & 0.48 & 0.16 \\
14 & 91 & 4.0 & 87 & 96\% & 0.46 & 0.14 \\
16 & 144 & 3.8 & 141 & 97\% & 0.45 & 0.12 \\
18 & 212 & 6.2 & 205 & 97\% & 0.43 & 0.15 \\
20 & 180 & 7.5 & 173 & 96\% & 0.37 & 0.15 \\
\end{longtable}

\textbf{Extended measurements} (C hash-table BFS, \(D = 10\)--\(28\), 10--30
samples per \(D\)). \emph{Data quality warning:} for \(D \geq 22\), the hash
table capacity limits the BFS tree size, causing semiprimes with larger \(S_0\)
to overflow and be excluded from the sample. The \(D \geq 22\) rows therefore
exhibit \textbf{downward survivor bias} in \(\alpha\) and should be treated as
lower bounds, not unbiased estimates:

\begin{longtable}[]{@{}rrrr@{}}
\toprule\noalign{}
\(D\) & \(n\) & \(\langle H_{\text{peak}} \rangle\) & \(\alpha = H/D\) \\
\midrule\noalign{}
\endhead
\bottomrule\noalign{}
\endlastfoot
10 & 30 & 4.98 & 0.498 \\
12 & 30 & 5.78 & 0.482 \\
14 & 30 & 6.92 & 0.494 \\
16 & 30 & 7.59 & 0.474 \\
18 & 20 & 7.94 & 0.441 \\
20 & 20 & 8.83 & 0.441 \\
22 & 15 & 8.34 & 0.379 \\
24 & 15 & 10.18 & 0.424 \\
26 & 10 & 10.35 & 0.398 \\
28 & 7 & 9.45 & 0.337 \\
\end{longtable}

Linear regression (\(D = 10\)--\(20\), highest-quality range):
\(H_{\text{peak}} \approx 0.38 D + 1.35\) (\(R^2 = 0.98\)). At \(D \geq 22\),
variance increases and \textbf{survivor bias} (semiprimes with higher
\(H_{\text{peak}}\) overflow the hash table and are excluded) causes the
observed slope to decrease.

\subsubsection{5.2 Scaling Summary}\label{scaling-summary}

\begin{longtable}[]{@{}lll@{}}
\toprule\noalign{}
Component & Slope (\(\alpha\)) & Asymptotic \\
\midrule\noalign{}
\endhead
\bottomrule\noalign{}
\endlastfoot
\(S_0\) (total) & \(0.38\)--\(0.50\) & Exponential \\
\(\tau\) (core) & \(\to 0\) & Sub-polynomial (Wigert--Ramanujan) \\
\(\Phi\) (phantom) & \(0.38\)--\(0.50\) & Exponential (dominates \(S_0\)) \\
\end{longtable}

\subsubsection{5.3 Growth Rate: Analytical Constraints and Empirical
Ambiguity}\label{growth-rate-analytical-constraints-and-empirical-ambiguity}

The analytical results constrain the phantom exponent
\(\alpha_\Phi = \lim_{D \to \infty} \log_2(\Phi_{\text{peak}})/D\):

\begin{itemize}
\tightlist
\item
  \textbf{Lower bound}: \(\alpha_\Phi > 0\) (Hensel degeneracy guarantees
  exponential modular aliasing).
\item
  \textbf{Upper bound}: \(\alpha_\Phi \leq 1\) (Proposition 3).
\item
  \textbf{Empirical range} (unbiased data, \(D = 10\)--\(20\)):
  \(\alpha \in [0.37, 0.50]\).
\end{itemize}

Within these bounds, three models fit the \(D = 10\)--\(20\) data with
comparable \(R^2\):

\begin{enumerate}
\def\labelenumi{\arabic{enumi}.}
\tightlist
\item
  \textbf{Linear}: \(H = \alpha D + b\), \(\alpha \approx 0.38\) (constant
  exponent)
\item
  \textbf{Sqrt correction}: \(H = D/2 - c \sqrt{D} + d\), \(\alpha \to 0.5\) as
  \(D \to \infty\)
\item
  \textbf{Log correction}: \(H = D/2 - c' \log_2 D + d'\), \(\alpha \to 0.5\) as
  \(D \to \infty\)
\end{enumerate}

The models diverge significantly only for \(D \geq 40\). At RSA-scale
(\(D = 2048\)), model 1 predicts \(H \approx 780\) while models 2--3 predict
\(H \approx 1000\). The biased data at \(D \geq 22\) cannot discriminate.
Whether \(\alpha_\Phi\) converges to \(1/2\) (the natural ``half-entropy''
threshold) or stabilizes below is the central open question of this paper.

\begin{center}\rule{0.5\linewidth}{0.5pt}\end{center}

\newpage

\subsection{\texorpdfstring{6. The \(p\)-adic Aliasing
Mechanism}{6. The p-adic Aliasing Mechanism}}\label{the-p-adic-aliasing-mechanism}

The BFS explores the 2-adic factorization tree: at step \(k\), it extends
partial solutions \(p_k \cdot q_k \equiv N \pmod{2^{k+1}}\), which is Hensel
lifting in \(\mathbb{Z}_2\).

In base 2, Hensel lifting is \emph{degenerate}: for any choice of bit
\(b_k(p)\), there exists a compatible \(b_k(q)\) maintaining the modular
congruence. This generates exponentially many modular solutions (the phantoms)
that satisfy local constraints but do not correspond to integer factorizations.

The real factorizations form a sub-polynomial subset \(\tau_{\text{odd}}(M_k)\)
embedded in the exponential phantom sea. The difficulty of factorization lies
not in the scarcity of true solutions (a semiprime always has exactly 2:
\((p,q)\) and \((q,p)\)) but in the exponential camouflage of the phantom halo.

\subsubsection{6.1 The Submarine Effect}\label{the-submarine-effect}

Experimentally, the true factorization path \((p, q)\) almost always passes
through intermediate BFS states with carry \(> 0\). The factors ``dive below''
the carry-zero surface and resurface only at the final bit. Enumerating
carry-zero states therefore does not directly yield a factoring algorithm.

\textbf{Remark (GCD heuristic).} As a tangential observation,
\(\gcd(N, |P_N(x_0)|)\) can recover a factor of \(N\) when \(x_0\) is a root of
\(P_N\) modulo \(p\) (exploiting \(P_N(2) = N \equiv 0 \pmod{p}\)). However,
binary polynomials have sparse root sets, yielding expected complexity
\(O(\sqrt{N} / \log N)\) --- not competitive with existing methods. We mention
this only as a curiosity arising from the polynomial representation; it is not a
contribution of this paper.

\begin{center}\rule{0.5\linewidth}{0.5pt}\end{center}

\newpage

\subsection{7. Implications and
Non-Implications}\label{implications-and-non-implications}

\subsubsection{What this establishes}\label{what-this-establishes}

\begin{itemize}
\tightlist
\item
  A structural decomposition of the carry-zero BFS tree: sub-polynomial core +
  exponential halo.
\item
  Identification of \(p\)-adic aliasing (Hensel degeneracy) as the mechanism
  generating phantom states.
\item
  A precise quantitative separation: phantoms dominate at \(> 93\%\) for
  \(D \geq 12\).
\end{itemize}

\subsubsection{What this does not imply}\label{what-this-does-not-imply}

\begin{itemize}
\tightlist
\item
  \textbf{No faster factoring algorithm}: the total BFS remains exponential; the
  carry-zero subset is structural, not algorithmic.
\item
  \textbf{No connection to RH}: these are combinatorial properties of binary
  arithmetic, not analytic properties of \(\zeta(s)\).
\end{itemize}

\begin{center}\rule{0.5\linewidth}{0.5pt}\end{center}

\newpage

\subsection{8. Open Questions}\label{open-questions}

\begin{enumerate}
\def\labelenumi{\arabic{enumi}.}
\tightlist
\item
  Exact growth rate of \(\Phi(k)\): Is
  \(\alpha_{\Phi} = \lim_{D \to \infty} \log_2(\Phi_{\text{peak}}) / D\)
  well-defined? What is its value?
\item
  \textbf{Carry-zero fraction}: What fraction of total BFS states have carry
  \(= 0\) at the peak? Does this fraction decay exponentially?
\item
  \textbf{Base dependence}: In bases \(b > 2\) where Hensel lifting is
  non-degenerate, how does the phantom fraction behave?
\end{enumerate}

\begin{center}\rule{0.5\linewidth}{0.5pt}\end{center}

\newpage

\subsection{9. Reproducibility}\label{reproducibility}

Key scripts in \texttt{experiments/}:

\begin{itemize}
\tightlist
\item
  \texttt{D01\_bfs\_entropy.c}: C hash-table BFS for \(D \leq 28\).
\item
  \texttt{D02\_hensel\_lifting.py}: Hensel lifting demonstration.
\item
  \texttt{D03\_phantom\_decomposition.py}: Phantom vs divisor decomposition and
  scaling exponents.
\item
  \texttt{D04\_phantom\_stats.py}: Detailed phantom statistics per \(N\).
\end{itemize}

Requirements: Python 3.8+, NumPy. C compiler (gcc/clang) for
\texttt{D01\_bfs\_entropy.c}.

\begin{center}\rule{0.5\linewidth}{0.5pt}\end{center}

\newpage

\subsection{References}\label{references}

\begin{enumerate}
\def\labelenumi{\arabic{enumi}.}
\tightlist
\item
  P. Diaconis, J. Fulman, ``Carries, Shuffling, and Symmetric Functions,''
  \emph{Adv. Appl. Math.} 43(2), 176--196, 2009.
\item
  S. Wigert, ``Sur l'ordre de grandeur du nombre des diviseurs d'un entier,''
  \emph{Arkiv for Matematik} 3(18), 1--9, 1907.
\item
  S. Ramanujan, ``Highly Composite Numbers,'' \emph{Proc. London Math. Soc.}
  2(1), 347--409, 1915.
\item
  {[}A{]} Companion paper: ``Spectral Theory of Carries in Positional
  Multiplication,'' this series.
\item
  {[}B{]} Companion paper: ``Carry Polynomials and the Euler Product: An
  Approximation Framework,'' this series.
\item
  {[}P1{]} Companion paper: ``π from Pure Arithmetic: A Spectral Phase
  Transition in the Binary Carry Bridge,'' this series.
\item
  {[}P2{]} Companion paper: ``The Sector Ratio in Binary Multiplication: From
  Markov Failure to Transcendence,'' this series.
\item
  {[}C{]} Companion paper: ``Eigenvalue Statistics of Carry Companion Matrices:
  A Markov-Driven GOE↔GUE Transition in Sparse Non-Hermitian Ensembles,'' this
  series.
\item
  {[}E{]} Companion paper: ``The Trace Anomaly of Binary Multiplication,'' this
  series.
\item
  {[}G{]} Companion paper: ``The Angular Uniqueness of Base 2 in Positional
  Multiplication,'' this series.
\item
  {[}H{]} Companion paper: ``Carry Polynomials and the Partial Euler Product: A
  Control Experiment,'' this series.
\end{enumerate}

\begin{center}\rule{0.5\linewidth}{0.5pt}\end{center}

\emph{CC BY 4.0. Code: MIT License.}

\end{document}
